\section{Conclusiones}
\label{sec:conclusiones}

\begin{table}[htbp]
    \centering
    \caption{Conclusiones del an\'{a}lisis de flexibilidad de la l\'{i}nea SIM-015.}
    \label{tab:conclusiones}
    \setcounter{itemcount}{0}
    \begin{tabularx}{\textwidth}{@{}NX@{}}
        \toprule
        \multicolumn{1}{c}{\textbf{Item}} & \textbf{Conclusion} \\
        \midrule
        & La tuber\'{i}a del filtro de fibras PP1 (SIM-015) cumple los criterios de esfuerzo de ASME B31.3-2016 en los nueve casos de carga evaluados: la evaluaci\'{o}n de cumplimiento de c\'{o}digo del modelo arroj\'{o} CODE COMPLIANCE EVALUATION PASSED. \\
        & El esfuerzo de c\'{o}digo m\'{a}ximo es \num{41819.1}~\si{kPa} en el nodo 80, caso 2 (Alt-SUS) W+P2, con una relaci\'{o}n de utilizaci\'{o}n de 36.3~\% frente al admisible de \num{115142.4}~\si{kPa}; la condici\'{o}n que gobierna el dise\~{n}o es la sostenida (peso m\'{a}s presi\'{o}n), no el rango de expansi\'{o}n t\'{e}rmica (ratio m\'{a}ximo de expansi\'{o}n 20.3~\% en el caso 8). \\
        & Los rangos de desplazamiento del sistema son moderados (m\'{a}ximo 10.0~\si{mm} en la componente $DX$ del caso 8, EXP) y las cargas en soportes son moderadas (carga vertical m\'{a}xima \num{5939}~\si{N} en el soporte del nodo 80 en operaci\'{o}n), por lo que la configuraci\'{o}n de soporter\'{i}a del isom\'{e}trico es adecuada. \\
        & La comparaci\'{o}n de las cargas en las conexiones con los equipos (anclas de los nodos 140, 320 y 510) con los admisibles de cada fabricante queda pendiente como verificaci\'{o}n externa, por no disponerse de dichos admisibles a la fecha. \\
        \bottomrule
    \end{tabularx}
\end{table}

En s\'{i}ntesis, la l\'{i}nea SIM-015 es t\'{e}cnicamente viable seg\'{u}n
ASME B31.3-2016 con una utilizaci\'{o}n m\'{a}xima de 36.3~\% del admisible.
No se identifican riesgos de sobre-esfuerzo, interferencias por
expansi\'{o}n ni sobrecarga de soportes, por lo que no se requieren acciones
correctivas sobre la configuraci\'{o}n analizada.
